\documentclass[conference]{IEEEtran}
\usepackage{cite}
\usepackage{amsmath,amssymb}
\usepackage{booktabs}
\usepackage{array}
\usepackage{multirow}
\usepackage{xcolor}
\usepackage{url}
\usepackage{microtype}

\newcolumntype{Y}{>{\raggedright\arraybackslash}p{0.22\textwidth}}
\newcolumntype{Z}{>{\raggedright\arraybackslash}p{0.35\textwidth}}

\begin{document}

\title{Why Do Humans Become Attached to Dogs and Cats?\\
BRIDGE: A Natural-Science Framework for Cross-Species Dyadic Affiliation}

\author{\IEEEauthorblockN{Design-Only Research Synthesis}}

\maketitle

\begin{abstract}
Humans often form durable, partner-specific bonds with dogs and cats, but a natural-science account must not convert a familiar everyday observation into a claim that all three species share identical emotions or mechanisms. This paper synthesizes evolutionary biology, comparative neuroethology, endocrinology, autonomic physiology, genetics, and animal welfare science. The evidence supports a constrained answer: domestication and cohabitation placed dogs and cats in human caregiving niches, while conserved mammalian systems for social recognition, reward, stress regulation, learning, and developmental plasticity can reinforce repeated contact with an individual animal. Dogs have the clearest direct evidence for a gaze-mediated, peripheral oxytocin-associated feedback process. Cats show partner-specific attachment-like behavior in Secure Base Test studies, but the available human-cat endocrine evidence is conditional and does not demonstrate a population-wide oxytocin increase. We introduce BRIDGE--Biological Routes to Interspecies Dyadic Guided Engagement--to separate evolutionary enablement, sensory recognition, integrated physiological state, dyadic regulation, genotype/development, and claim calibration. A preregistered multimodal design is specified for 180 established dog-owner and cat-owner dyads. It compares owners, familiar nonowners, barrier, and quiet-control conditions while measuring behavior, autonomic dynamics, validated peripheral assays, and development. The proposal makes no new empirical claim. Its central contribution is a falsifiable route to determine which mechanisms are shared, which are species-specific, and which popular explanations should be rejected.
\end{abstract}

\begin{IEEEkeywords}
human-animal interaction, dog, cat, domestication, oxytocin, autonomic physiology, comparative neuroethology, attachment behavior, animal welfare
\end{IEEEkeywords}

\section{Introduction}

Why humans become attached to dogs and cats is scientifically tractable only after the question is narrowed. ``Attachment'' is not used here as a philosophical category, a cultural preference, or proof of matching subjective experience. It denotes a measurable pattern of partner-specific approach, proximity maintenance, separation and reunion behavior, and physiological regulation in a stable dyad. The target is therefore an organismal question: why does a familiar dog or cat become a reliable stimulus for human caregiving and affiliation, and why can that particular human become a salient social partner for the animal?

The answer cannot be a single hormone or a single domestication event. Dogs and cats entered human environments through different ecological histories. Dog domestication included selection on behavior and communication in a human social niche, with genomic evidence of adaptation to human-associated diet and ecology \cite{axelsson2013}. Cats most likely spread through commensal association with agricultural settlements and later transport by people, while retaining marked flexibility in social organization \cite{driscoll2007,ottoni2017}. The routes converge on repeated contact with humans but do not imply the same sensory biases, developmental dependencies, or neuroendocrine responses.

The strongest dog-specific experimental evidence is the report by Nagasawa \emph{et al.} that dog-owner mutual gaze was associated with urinary oxytocin changes and that intranasal oxytocin administered to dogs increased gaze in their experimental setting \cite{nagasawa2015}. This supports a candidate interspecies loop, not a universal diagnostic of love or a measurement of central brain oxytocin. For cats, Vitale \emph{et al.} reported Secure Base Test behavior consistent with distinct attachment styles toward human caregivers in kittens and adult cats \cite{vitale2019}. Johnson \emph{et al.}, in contrast, found no group-average salivary oxytocin change in 30 women during interaction with their pet cats compared with reading, although selected interaction behaviors covaried with hormone change \cite{johnson2021}. These results make a useful scientific point: dogs and cats should not be ranked on one anthropomorphic scale, nor should a dog result be automatically transferred to cats.

This design-only paper organizes the literature and proposes a protocol. Section II establishes what present evidence permits. Section III defines the BRIDGE framework. Sections IV and V turn the framework into preregistered measurements and models. Section VI specifies what conclusions would be justified or ruled out. No new samples, hormone values, genetic effects, behavioral effect sizes, health outcomes, or causal neural findings are claimed.

\section{What Existing Evidence Establishes}

\subsection{Evolutionary enablement is not an individual-level explanation}

Domestication changes the distribution of traits on which learning acts; it does not guarantee a particular dyad will form a bond. Dogs are especially informative because socialized dogs can differ from hand-reared wolves in human-directed behavior, and a recent companion pig-dog comparison found attachment-specific owner responses in dogs but not pigs under its protocol \cite{topal2005,gabor2024}. The comparative result rejects a simple explanation in which general domestication plus exposure to humans is sufficient. It does not isolate a single gene, prove a dog-only capacity, or establish a mechanism in cats.

For cats, a commensal origin and flexible social ecology predict larger individual and context dependence. Feral and free-ranging cats may live alone or in groups as resource distribution changes. In a household, access to food, predictable resting sites, tactile contact, and low-threat handling can repeatedly pair an individual human with safety and reward. The correct inference is an ecological opportunity for affiliation, not an assertion that cats were selected for the same human-oriented gaze system as dogs.

Behavioral definitions provide a bridge across species. In a Secure Base Test or related design, an attachment-relevant response includes different behavior toward a caregiver and a stranger, changes during separation, and organized reunion behavior. Dog-owner work has found owner-directed proximity and physiological covariation during a modified Strange Situation Test \cite{ryan2019,nagasawa2009}; dog secure-base behavior has also been reported under owner-presence manipulations \cite{horn2013}. Cat Secure Base Test evidence makes equivalent behavioral questions testable for cats, provided the ethogram respects feline sensory and motor behavior \cite{vitale2019}. A human observer's impression of affection is neither necessary nor sufficient.

\subsection{Recognition, learning, and social buffering}

A viable biological account requires individual recognition. Dogs use human visual and gestural signals at levels that made them a central model for comparative social cognition \cite{hare2005,miklosi2013,udell2016}. Repeated exposure lets a particular person's odor, voice, movement, face, touch pattern, and temporal routine become predictive cues. Predictive learning can increase voluntary approach without requiring any claim about verbal concepts. Awake, unrestrained canine functional MRI also establishes the technical feasibility of future circuit-level hypothesis tests, but it does not itself explain companion attachment \cite{berns2012}. The same logic applies to cats, although the weight of cues may differ: voluntary approach initiation, contact choice, posture, vocalization, odor, and touch are likely more diagnostic than sustained mutual gaze alone.

The proposed mechanism is not merely reinforcement by food. Food expectation can produce approach to any feeder; attachment requires a contrast in which the owner explains behavior beyond a familiar nonowner, equal food expectation, shared activity, and novelty. Positive human-animal relationships are commonly indexed by voluntary approach, proximity, anticipation, relaxation, and choice, but these indicators remain sensitive to prior handling and individual predisposition \cite{rault2020}. Thus, familiar-partner recognition is a necessary measurement target rather than a narrative embellishment.

\subsection{Peripheral peptides and autonomic state}

Oxytocin is relevant because it has conserved roles in mammalian parturition, lactation, parental care, social recognition, and stress regulation. It is not a scalar ``bonding chemical.'' Its action depends on receptor distribution, sex steroids, developmental stage, context, species, and neural circuit. In dogs, the gaze study provides a biologically coherent candidate loop: dog-to-owner gaze predicted urinary oxytocin changes and dog oxytocin administration affected gaze under controlled conditions \cite{nagasawa2015}. Other dog-owner studies report attachment-related behavior, cortisol or chromogranin A measures, and hormonal covariation, but the direction and magnitude depend on the challenge and assay \cite{ryan2019,handlin2015}.

Cat evidence is deliberately more cautious. The 2021 cat study found behavioral correlates of women's salivary oxytocin change but not the expected group-level increase relative to a reading control \cite{johnson2021}. That result does not mean cat relationships lack biology; it means a simple ``interaction raises oxytocin'' explanation was not supported in that sample and protocol. It also demonstrates why saliva, urine, and plasma results must be reported as peripheral associations. They cannot directly identify synaptic release, receptor occupancy, or a common neural state across species.

Autonomic measures offer complementary temporal resolution. Heart rate, heart-rate variability (HRV), electrodermal activity, movement, and cortisol can reveal whether interaction lowers or raises arousal, but their interpretation remains context dependent. A coincident rise in HRV or fall in cortisol can be caused by rest, reduced movement, room familiarity, circadian timing, or experiment completion. A dyadic mechanism requires controls for these common causes and time-lagged analysis rather than a visually appealing correlation.

\begin{table*}[t]
\caption{Evidence levels for a natural-science account of human attachment to dogs and cats.}
\label{tab:claim-ladder}
\centering
\footnotesize
\begin{tabular}{p{0.14\textwidth}p{0.22\textwidth}p{0.24\textwidth}p{0.19\textwidth}}
\toprule
\textbf{Claim level} & \textbf{Permitted statement} & \textbf{Minimum evidence} & \textbf{Alternative that must remain visible} \\
\midrule
Ecological enablement & A species can exploit a human-associated niche. & Domestication history, comparative ecology, or genomic context. & Individual attachment is not implied. \\
Partner-specific behavior & A dyad behaves differently with the owner than with a matched person. & Counterbalanced owner, familiar-nonowner, and control trials with a preregistered ethogram. & Food anticipation, novelty, training, handler cueing, or general sociability. \\
Peripheral physiological association & A specified peripheral measure covaries with interaction features. & Validated collection, assay controls, time-locking, and reporting of failed samples. & Peripheral measure is not central neural signaling or causation. \\
Candidate dyadic regulation & One partner's state improves prediction of the other's later state. & Synchronized time series beating shared-room, movement, and schedule null models. & Shared environment and coupled activity. \\
Molecular or neural mechanism & A receptor, circuit, or developmental process mediates the effect. & Targeted perturbation or independently replicated longitudinal evidence. & Population stratification, off-target effects, and species-specific circuitry. \\
\bottomrule
\end{tabular}
\end{table*}

\section{BRIDGE: A Causal Organization of the Problem}

BRIDGE--Biological Routes to Interspecies Dyadic Guided Engagement--is a measurement architecture rather than an additional theory of emotion. Figure \ref{fig:bridge} orders the evidence so that a failure at one stage narrows a claim instead of erasing a lower-level observation. Background selection (B) explains why a species may be prepared for a human niche. Recognition (R) tests whether a particular human is detected as a partner. Integrated state (I) tests behavior with autonomic and endocrine measures. Dyad regulation (D) asks whether time-resolved coupling exceeds shared context. Genotype and development (G) model why effects vary. Evidence calibration (E) binds every conclusion to controls, replication, and the claim ladder in Table \ref{tab:claim-ladder}.

\begin{figure*}[t]
\centering
\footnotesize
\fbox{\begin{minipage}{0.94\textwidth}\centering
\fbox{\begin{minipage}{0.16\textwidth}\centering\textbf{B: Background}\\Domestication, ecology, development\end{minipage}}
\hfill $\Rightarrow$ \hfill
\fbox{\begin{minipage}{0.16\textwidth}\centering\textbf{R: Recognition}\\Owner cues versus matched person\end{minipage}}
\hfill $\Rightarrow$ \hfill
\fbox{\begin{minipage}{0.16\textwidth}\centering\textbf{I: Integrated state}\\Behavior, autonomic, peripheral assays\end{minipage}}
\hfill $\Rightarrow$ \hfill
\fbox{\begin{minipage}{0.16\textwidth}\centering\textbf{D: Dyad regulation}\\Time-lagged prediction beyond context\end{minipage}}
\par\vspace{6pt}
\fbox{\begin{minipage}{0.21\textwidth}\centering\textbf{G: Genotype and development}\\Moderation, not a direct causal claim\end{minipage}}
\hfill $\Rightarrow$ \hfill
\fbox{\begin{minipage}{0.25\textwidth}\centering\textbf{E: Evidence calibration}\\Replicate, compare nulls, report the highest supported level\end{minipage}}
\end{minipage}}
\caption{BRIDGE separates enabling evolutionary context from individual recognition, measured physiology, dyadic regulation, moderation, and calibrated inference. Dog and cat analyses remain separate until a pre-specified cross-species comparison is justified.}
\label{fig:bridge}
\end{figure*}

\subsection{A preregistered behavioral endpoint}

For each species $s$ and dyad $i$, the primary endpoint is an affiliation index constructed only from measures feasible and welfare-compatible for that species:
\begin{equation}
B_{is}=\sum_{k=1}^{K_s}w_{ks}z_{iks}-\sum_{m=1}^{M_s}v_{ms}z_{ims}.
\label{eq:affiliation-index}
\end{equation}
In \eqref{eq:affiliation-index}, $z_{iks}$ are favorable standardized measures such as voluntary approach, partner-directed contact, or organized reunion behavior; $z_{ims}$ are withdrawal, persistent avoidance, or distress measures; and the weights $w_{ks}$ and $v_{ms}$ are locked before enrollment. Species-specific components avoid falsely treating feline gaze duration as equivalent to canine gaze duration. The endpoint is not a claim of subjective intensity.

The owner effect is the within-dyad contrast
\begin{equation}
\Delta B_{is}=B_{is}^{\mathrm{owner}}-B_{is}^{\mathrm{familiar\mbox{-}nonowner}}.
\label{eq:owner-contrast}
\end{equation}
Equation \eqref{eq:owner-contrast} eliminates a broad class of explanations based on any human presence. A positive $\Delta B_{is}$ must still be compared with the barrier and quiet-control conditions before it is called partner specific.

\subsection{Predictions that can fail}

The first prediction is that owner interaction will produce a positive mean of \eqref{eq:owner-contrast} in dogs and cats. The second predicts a species-by-cue interaction: mutual gaze will contribute more to dog prediction, while voluntary approach and tactile choice will contribute more to cat prediction. The third predicts that peripheral oxytocin and cortisol effects are conditional on baseline state, assay matrix, and interaction quality, not uniformly positive. The fourth predicts that any dyadic synchrony will outperform models based only on common movement and shared schedule. The fifth treats genotype and early socialization as moderators; their associations are exploratory until reproduced independently.

\section{A Multimodal, Welfare-First Test}

\subsection{Sampling and conditions}

The confirmatory cohort contains 90 established dog-owner and 90 established cat-owner dyads, each with a familiar nonowner when available. A separate longitudinal substudy follows 40 newly formed dyads from the first week after adoption through month six; its estimates cannot be pooled with established dyads. Enrolling both familiar and nonowner conditions is central: a pet that approaches anyone is socially responsive, whereas a pet that shows an owner-specific contrast satisfies a stronger biological observation.

Each established dyad completes four counterbalanced 12-minute sessions at the same circadian window on separate days: owner free interaction, familiar-nonowner free interaction, owner behind a transparent barrier without touch, and quiet rest with a neutral object. Food, novel toys, verbal prompts, room temperature, session order, movement space, and collection timing are fixed or recorded. The animal may disengage at any time; there is no deprivation, restraint, deliberate separation beyond the minimal protocol, or intranasal hormone administration. These choices make a negative result interpretable as well as safe.

\begin{table}[t]
\caption{Pre-specified modalities and the claim each can support.}
\label{tab:modalities}
\centering
\footnotesize
\begin{tabular}{p{0.21\columnwidth}p{0.28\columnwidth}p{0.30\columnwidth}}
\toprule
\textbf{Modality} & \textbf{Recorded variable} & \textbf{Permitted inference} \\
\midrule
Blinded video coding & gaze/head orientation, voluntary approach, contact, withdrawal, vocalization, movement & owner-specific behavior after matched controls \\
Autonomic sensors & HRV, heart rate, electrodermal activity where feasible, activity & time-resolved arousal and candidate co-regulation \\
Peripheral assays & saliva/urine cortisol and validated oxytocin protocol & matrix-specific association, not central neural release \\
Development and health & age, sex, neuter status, medication, sleep, illness, socialization, caregiving duration & covariate adjustment and effect moderation \\
Optional SNP panel & ancestry-aware OXTR and related variants & exploratory association requiring replication \\
\bottomrule
\end{tabular}
\end{table}

\subsection{Physiology and assay discipline}

For a physiological quantity $X$ measured at baseline $t_0$ and after condition $c$ at $t_1$, the pre-specified response is
\begin{equation}
\Delta X_{isc}=X_{isc}(t_1)-X_{isc}(t_0).
\label{eq:phys-change}
\end{equation}
Equation \eqref{eq:phys-change} is reported for each assay matrix rather than converting saliva, urine, and plasma into a single unvalidated score. Samples with time, volume, storage, or assay-quality failures remain visible in a flow record. Assay precision is summarized by the coefficient of variation
\begin{equation}
\mathrm{CV}_{q}=\frac{\sigma_q}{\mu_q},
\label{eq:assay-cv}
\end{equation}
where $\mu_q$ and $\sigma_q$ are the mean and standard deviation of quality-control replicates for assay batch $q$. Equation \eqref{eq:assay-cv} is a quality criterion, not a biological outcome.

Behavioral coders are blinded to condition labels, and an independent reliability subset is coded twice. Reliability is stated as an intraclass correlation:
\begin{equation}
\mathrm{ICC}=\frac{\sigma^2_{\mathrm{dyad}}}{\sigma^2_{\mathrm{dyad}}+\sigma^2_{\mathrm{coder}}+\sigma^2_{\epsilon}}.
\label{eq:icc}
\end{equation}
In \eqref{eq:icc}, the variance components represent dyad, coder, and residual contributions. A low value triggers endpoint review before biological interpretation; it is not repaired by post hoc recoding choices.

\subsection{Primary hierarchical model}

The primary analysis models species separately and then tests a declared interaction:
\begin{align}
\Delta B_{is}={}&\beta_{0s}+\beta_{1s}\mathrm{Owner}_{is}+\beta_{2s}\mathrm{Cue}_{is}
 +\beta_{3s}\mathrm{Order}_{is}\nonumber\\
&+\boldsymbol{\gamma}_{s}^{\mathsf{T}}\mathbf{C}_{is}+u_{i} +\epsilon_{is}.
\label{eq:mixed-model}
\end{align}
In \eqref{eq:mixed-model}, $\mathbf{C}_{is}$ includes health, development, reproductive status, time of day, and movement covariates; $u_i$ is a dyad random intercept; and $\epsilon_{is}$ is residual variation. The coefficient $\beta_{1s}$ estimates the owner contrast within species. A cross-species difference is tested only after model diagnostics demonstrate that the relevant behavioral measures are comparable enough for that narrow contrast.

Autonomic synchrony is initially a descriptive, lag-free quantity:
\begin{equation}
r_{is}^{(0)}=\mathrm{cor}\left(H^{\mathrm{human}}_{is}(t),H^{\mathrm{animal}}_{is}(t)\right),
\label{eq:zero-lag}
\end{equation}
where $H(t)$ is an aligned HRV or heart-rate-derived series. Equation \eqref{eq:zero-lag} cannot distinguish shared movement from regulation. The candidate directional model is therefore
\begin{align}
H^{\mathrm{animal}}_{is}(t+\ell)={}&a_0+a_1H^{\mathrm{animal}}_{is}(t)+a_2H^{\mathrm{human}}_{is}(t)\nonumber\\
&+\mathbf{a}_{3}^{\mathsf{T}}\mathbf{W}_{is}(t)+e_{is}(t).
\label{eq:cross-lag}
\end{align}
In \eqref{eq:cross-lag}, $\ell$ is a pre-specified lag and $\mathbf{W}(t)$ contains movement, room, condition, and schedule terms. The apparent $\mathbf{a}_3$ formatting in the equation is a scalar coefficient vector applied to the covariates; it is tested against a shared-exposure null. Evidence is reported as a prediction improvement, not as proof that one organism controls the other.

\section{Model Competition, Causal Boundaries, and Reporting}

\subsection{Competing explanations}

Three model families must compete. The \emph{partner-recognition} model predicts that owner identity explains affiliation after familiar-person controls. The \emph{general-reward} model predicts similar approach to any person associated with pleasant handling or food expectation. The \emph{shared-context} model predicts physiological correlation from the same room, movement pattern, or sample schedule. None of these models needs to deny that the other process exists; the question is which improves prediction of a declared outcome.

For a contrast in which $Y_i(1)$ is the outcome under owner interaction and $Y_i(0)$ the matched familiar-nonowner outcome, the target causal quantity is
\begin{equation}
\tau_s=\mathbb{E}\left[Y_i(1)-Y_i(0)\mid s\right].
\label{eq:causal-estimand}
\end{equation}
Equation \eqref{eq:causal-estimand} is identifiable only to the extent that counterbalancing, measured covariates, compliance, and missingness assumptions hold. It does not identify the neural mediator. In particular, a peripheral peptide association is not the mediator estimate in \eqref{eq:causal-estimand}.

Out-of-sample performance guards against a flexible story fitting every dyad. For held-out observations $j=1,\ldots,J$, the Brier score for a preregistered binary behavioral event is
\begin{equation}
\mathrm{BS}=\frac{1}{J}\sum_{j=1}^{J}\left(p_j-y_j\right)^2,
\label{eq:brier}
\end{equation}
where $p_j$ is the model probability and $y_j$ is the observed event. Equation \eqref{eq:brier} favors calibrated prediction rather than an appealing but untestable mechanism. A dyadic-regulation claim advances only if its held-out score improves over the owner-label-only and shared-context models.

The evidence grade is a transparent composite rather than an automatic truth machine:
\begin{equation}
E=I_{\mathrm{pre}}I_{\mathrm{quality}}I_{\mathrm{control}}I_{\mathrm{replication}}.
\label{eq:evidence-grade}
\end{equation}
In \eqref{eq:evidence-grade}, each indicator is one only when preregistration, assay/data quality, relevant controls, and replication criteria are met, respectively. If any required element is zero, $E=0$ and the associated mechanistic claim is downgraded. The product deliberately prevents a large sample from compensating for a missing control.

\begin{table*}[t]
\caption{Decision rules that stop overinterpretation.}
\label{tab:decisions}
\centering
\footnotesize
\begin{tabular}{p{0.22\textwidth}p{0.30\textwidth}p{0.37\textwidth}}
\toprule
\textbf{Observed pattern} & \textbf{What may be reported} & \textbf{What must not be reported} \\
\midrule
Owner contrast in behavior only & Partner-specific behavioral affiliation under the tested conditions. & A shared endocrine or neural mechanism. \\
Behavior plus peripheral association & A matrix- and protocol-specific physiological association. & That peripheral oxytocin is brain oxytocin or is the cause of attachment. \\
Zero-lag HRV correlation & Physiological covariation during a session. & Bidirectional regulation or emotional contagion. \\
Cross-lag model beats shared-context null on held-out data & Candidate dyadic temporal coupling. & Direct biological control without an intervention. \\
Exploratory OXTR association & Replication target and possible moderation. & A gene for pet attachment or a species-wide causal effect. \\
Dog result differs from cat result & Species- and protocol-specific divergence. & A rank of affection, intelligence, or capacity for experience. \\
\bottomrule
\end{tabular}
\end{table*}

\subsection{Why dog and cat estimates must remain distinct}

Dogs may show a strong owner-gaze pathway because their domestication and training history made gaze a high-value social signal. Cats may form an owner-specific relationship through different combinations of proximity, touch, vocal or olfactory signals, and environmental predictability. This is a testable divergence, not a stereotype. The cat endocrine result already shows why pooling can obscure the science: a null group-average oxytocin change can coexist with behavior-dependent associations \cite{johnson2021}. Conversely, the dog oxytocin-gaze result should be reported with its experimental context, peripheral assay limits, and wolf comparison rather than used as an all-purpose explanation \cite{nagasawa2015}.

Genetic evidence deserves equal restraint. OXTR variation in dogs and owners has been associated with components of dog-owner attachment in one comparative study \cite{kovacs2018}. Such a result motivates preregistered replication with ancestry, breed, relatedness, and environment controls. It does not identify a universal ``attachment gene.'' Early developmental experience is also biologically consequential in dogs \cite{dietz2018}; a future cat analysis must establish equally appropriate developmental measures rather than borrowing dog variables by default.

\section{Discussion}

The natural-science answer is consequently plural but precise. Humans become attached to dogs and cats when a domesticated animal repeatedly becomes an individual, multisensory predictor of safety, care, and reward, while the human receives reliable behavioral cues that recruit conserved caregiving, reward, and stress-regulation systems. Co-residence and history supply repeated learning opportunities; domestication supplies species-specific predispositions; development and health tune the response; and each dyad develops within physical constraints of sensory ecology and welfare. These are biological propositions because each can be measured, perturbed in ethically permissible ways, or falsified against alternatives.

The evidence is strongest for dog-human gaze and its peripheral oxytocin-associated feedback under a particular experimental design. It is strong for cat-caregiver behavioral attachment capacity, but current cat endocrine data do not justify a direct transfer of the dog loop. The correct explanation for both animals is therefore neither ``the same as human family bonds'' nor ``mere feeding.'' It is a family of cross-species dyadic mechanisms that must be resolved one level at a time.

BRIDGE makes that resolution auditable. It prevents ecological history from being misreported as individual causation, prevents peripheral peptide values from being misreported as brain states, and prevents physiological correlation from being misreported as regulation. It also identifies an empirical path forward: collect synchronized, species-appropriate behavior and physiology in owner, familiar-nonowner, barrier, and neutral-control conditions; preregister the claims; and require held-out success before promoting a mechanism. Such work can improve basic understanding while protecting animal welfare and preserving scientific humility about experience.

\appendices
\section{Preregistration and Welfare Record}

Before recruitment, the study locks the species-specific ethogram, primary endpoint in \eqref{eq:affiliation-index}, counterbalancing schedule, sample exclusions, assay SOPs, missing-data policy, model in \eqref{eq:mixed-model}, null models for \eqref{eq:cross-lag}, and the evidence gates in \eqref{eq:evidence-grade}. Any later analysis is labeled exploratory. A veterinarian and a human-research oversight body review the protocol. The study records adverse welfare events, voluntary disengagement, and sample failures rather than silently discarding them.

\section{Mechanistic Time Scales}

The same observed owner approach can arise at different biological time scales, and the experiment should not confuse them. At the scale of seconds, a dog or cat can orient to a familiar odor, voice, face, posture, or movement and choose to approach or withdraw. These events are appropriate for time-stamped video and autonomic analysis. At the scale of minutes, touch, gaze, vocal exchange, rest, or novelty can change movement, heart rate, HRV, and measurable peripheral analytes. Such changes are condition-specific physiological observations, not proof of a long-lived bond. The dog gaze experiment belongs chiefly at this minutes scale, although its domestication interpretation connects it to a longer history \cite{nagasawa2015}.

At the scale of days to months, repeated predictability can alter learned approach, handling tolerance, and the probability that a particular human functions as a safe, familiar partner. This is where the longitudinal adoption substudy is essential: it records how a dyad changes instead of inferring developmental history from a single established relationship. At the scale of development and generations, early social exposure, maternal care, selection, population structure, and domestication shape the range of responses available to an individual. The cat and dog domestication records explain this background but cannot replace a within-dyad contrast \cite{axelsson2013,driscoll2007,ottoni2017}.

The time-scale separation dictates the data architecture. Short autonomic windows are analyzed around observable cue events. Peripheral samples are interpreted only at their declared collection interval. Behavioral change across visits is modeled as a trajectory, not as a delayed hormone value. Genetic and developmental variables are baseline moderators, not rapidly changing mediators. A claim can move from one time scale to another only with an explicit measurement bridge. For example, a genetic association may motivate an early-life or receptor-expression study, but it cannot by itself explain a 30-second gaze event.

\begin{table}[t]
\caption{Time-scale map for BRIDGE measurements.}
\label{tab:timescales}
\centering
\footnotesize
\begin{tabular}{p{0.18\columnwidth}p{0.27\columnwidth}p{0.31\columnwidth}}
\toprule
\textbf{Scale} & \textbf{Observable} & \textbf{Permitted conclusion} \\
\midrule
Seconds & cue orientation, approach, withdrawal, event-locked HRV & sensory/behavioral response to the tested partner \\
Minutes & interaction sequence, movement, peripheral assay change & protocol-specific physiological association \\
Weeks onward & repeated owner contrast, socialization trajectory & developmental change in partner-specific behavior \\
Generations & domestication history, genomic context & enabling species background, not an individual mechanism \\
\bottomrule
\end{tabular}
\end{table}

\section{Evidence and Deviation Ledger}

The accompanying ledger labels every sentence as one of: publisher-verified evidence, literature synthesis, design proposal, or limitation. It records the source population, species, assay matrix, behavioral task, and inference boundary. The ledger rejects a claim of common dog-cat mechanism unless an explicitly comparable measure has passed reliability and control criteria in both species. Data release preserves reproducible behavioral code and assay metadata while excluding re-identifiable human genotype, household location, and animal medical records.

\section{Falsification Matrix and Recovery Analysis}

The protocol is useful only if it can reject attractive explanations. A strong owner effect with no difference between the owner and familiar-nonowner condition favors general sociability or a broadly rewarding human interaction over partner specificity. An owner effect present only when food is available favors resource prediction. An autonomic association that disappears after movement, room, and schedule adjustment favors shared exposure. A peptide association that changes sign across assay matrices or fails batch-quality criteria is an assay or protocol result, not evidence for a conserved social mechanism. A dog effect that is absent in cats is an empirical species difference, not a deficit in cats or a ranking of the two relationships.

Before enrollment, a simulation suite generates synthetic dyads with known combinations of partner-recognition, food-expectation, shared-movement, and measurement-error effects. The suite tests whether the planned model in \eqref{eq:mixed-model} can recover the owner coefficient without promoting a shared-context signal, and whether the cross-lag model in \eqref{eq:cross-lag} improves the held-out score in \eqref{eq:brier} only when a directional input was placed in the simulation. Recovery is evaluated separately for dogs and cats because their ethograms and sampling feasibility differ. If the procedure cannot recover an injected partner effect or routinely creates a false dyadic-regulation conclusion, the relevant confirmatory endpoint is removed or redesigned before animals are enrolled.

The simulation also varies missing data. Pets may refuse saliva collection, move out of a sensor field, or terminate a session; these events are biologically and ethically meaningful. A valid analysis cannot assume that missingness is random merely because the data file is incomplete. The preregistration therefore records each missingness mechanism, compares complete-case and sensitivity analyses, and reports the fraction of dyads contributing to every model. No imputation procedure is allowed to fabricate an endocrine response after a failed sample collection.

\begin{table}[t]
\caption{Falsification cases in the BRIDGE protocol.}
\label{tab:falsification}
\centering
\footnotesize
\begin{tabular}{p{0.22\columnwidth}p{0.25\columnwidth}p{0.26\columnwidth}}
\toprule
\textbf{Pattern} & \textbf{Favored explanation} & \textbf{Required report language} \\
\midrule
Owner and familiar person yield equal $\Delta B$ & General human-directed sociability & No partner-specific attachment conclusion \\
Effect is food-dependent only & Resource prediction & Food-associated approach, not dyadic regulation \\
Synchrony vanishes with activity covariate & Shared movement/context & No evidence for time-lagged co-regulation \\
Peripheral assay fails QC or replication & Assay or protocol instability & No peptide mechanism conclusion \\
Dog/cat estimates differ & Species-specific cue architecture & Divergent result under the tested protocol \\
\bottomrule
\end{tabular}
\end{table}

\section{Signal Processing and Molecular Boundaries}

The study treats a recorded signal as a measurement process, not as a direct window into motive. Video is time-stamped before coding, and raw sensor data retain clock offsets, dropped packets, and calibration records. HRV requires a species-appropriate interbeat-interval pipeline; a device that is acceptable for a resting dog may not be valid for a moving cat. Hormone samples require collection-to-freezer time, storage cycle, extraction method, assay batch, duplicate agreement, and lower-limit-of-quantification status. A sample label is never silently converted into a value because a model expects complete rows.

Molecular measures are particularly easy to overinterpret. Optional genotyping is limited to a predeclared moderation analysis, performed only with consent and an ancestry-aware plan. It is not a search for a gene that makes humans love pets. OXTR variants may influence receptor regulation, linkage patterns, or nothing measurable in a given population; breed structure in dogs and demographic structure in people can create associations even when no dyadic mechanism is involved. The safest progression is discovery in one cohort, preregistered replication in another, and only then a species-appropriate functional study. That sequence is more informative than adding a large genomic panel to an underpowered behavioral experiment.

Immune and microbiome variables are similarly kept outside the primary claim. Cohabitation can alter microbial exposure and immune environment, but this does not establish that microbes cause attachment or that attachment causes a health benefit. A later amendment may collect such samples only when it specifies a measurement model, plausible confounders, and a welfare/consent plan. Their omission from the confirmatory endpoint protects the central behavioral and physiological contrasts from uncontrolled multiple testing.

\section{Interpretation Record for Negative and Mixed Results}

The most valuable outcome may be a mixed result. Suppose dogs show a reproducible owner-gaze association but cats show an owner-proximity effect with no group-level oxytocin change. BRIDGE would support a shared high-level proposition--partner-specific affiliation can be biologically organized in both species--while rejecting a shared low-level oxytocin-gaze mechanism. Suppose both species show behavior but neither physiology nor time-lagged prediction survives controls. The correct conclusion would be behavioral partner specificity under the tested conditions, with no evidence yet for a common physiological pathway. Conversely, no owner contrast in a healthy, sufficiently powered cohort would challenge the proposed endpoint or sample definition; it would not show that every owner-pet relationship is biologically empty.

This record protects welfare as well as inference. It makes voluntary disengagement a scientifically visible outcome, prohibits stressful manipulation to force attachment-like behavior, and separates a pet's approach from a human's desire for a bond. A natural-science explanation should expand what can be tested while leaving unmeasured experience unclaimed.

\section{Cross-Species Measurement Card}

Every released analysis receives a compact measurement card. The card names the species, life stage, sex and reproductive-status handling, health exclusions, recruitment route, duration of cohabitation, condition order, room properties, and all departures from the planned session. It then lists the exact behavioral variables used in $B_{is}$ from \eqref{eq:affiliation-index}, the coder blinding status and reliability from \eqref{eq:icc}, sensor model and sampling rate, assay matrix, batch identifiers, and the number of valid observations at each time point. This level of documentation is necessary because two studies can both use the word ``interaction'' while observing entirely different biology.

The card also distinguishes a construct from its proxy. Mutual gaze is an observable visual event, not a direct measurement of affiliative neural activity. Salivary oxytocin is a peripheral analyte, not a direct measurement of hypothalamic release. HRV is a quantified property of a cardiac time series, not a direct scale of calmness. The distinction avoids a common chain of errors: replacing a controlled observation with a convenient label, then treating the label as the mechanism. BRIDGE requires the original variable, its sensor or assay, and its uncertainty to travel together in every figure and model.

Equivalence is assessed before comparison. A cat's voluntary approach and a dog's mutual gaze may both be informative but need not occupy the same coordinate in a latent attachment construct. The primary results therefore begin with within-species owner contrasts and condition effects. Only a prespecified measurement-invariance or prediction-transfer analysis may support a pooled statement. If such an analysis fails, species-specific estimates are the result, not a methodological failure. This preserves the possibility that different sensory pathways converge on similar dyadic stability.

Finally, the card records the direction of uncertainty. A failed assay, an ambiguous video segment, or an animal that chooses not to engage makes a molecular or dyadic conclusion less certain; it does not justify discarding the event or reclassifying it after outcomes are known. Combined with the evidence grade in \eqref{eq:evidence-grade}, this record turns reproducibility from a final editorial check into a biological safeguard. It helps future studies determine whether a discrepancy reflects species, protocol, measurement, or a genuine limit of the proposed mechanism.

\begin{thebibliography}{99}

\bibitem{nagasawa2015}
M. Nagasawa \emph{et al.}, ``Oxytocin-gaze positive loop and the coevolution of human-dog bonds,'' \emph{Science}, vol. 348, no. 6232, pp. 333--336, 2015, doi: 10.1126/science.1261022.

\bibitem{vitale2019}
K. R. Vitale, A. C. Behnke, and M. A. R. Udell, ``Attachment bonds between domestic cats and humans,'' \emph{Current Biology}, vol. 29, no. 18, pp. R864--R865, 2019, doi: 10.1016/j.cub.2019.08.036.

\bibitem{johnson2021}
E. A. Johnson, A. Portillo, N. E. Bennett, and P. B. Gray, ``Exploring women's oxytocin responses to interactions with their pet cats,'' \emph{PeerJ}, vol. 9, e12393, 2021, doi: 10.7717/peerj.12393.

\bibitem{gabor2024}
A. Gabor, P. Perez Fraga, M. Gacsi, L. Gerencser, and A. Andics, ``Domestication and exposure to human social stimuli are not sufficient to trigger attachment to humans: a companion pig-dog comparative study,'' \emph{Scientific Reports}, vol. 14, art. 14058, 2024, doi: 10.1038/s41598-024-63529-3.

\bibitem{ryan2019}
M. G. Ryan, A. E. Storey, R. E. Anderson, and C. J. Walsh, ``Physiological indicators of attachment in domestic dogs and their owners in the Strange Situation Test,'' \emph{Frontiers in Behavioral Neuroscience}, vol. 13, art. 162, 2019, doi: 10.3389/fnbeh.2019.00162.

\bibitem{kovacs2018}
K. Kovacs \emph{et al.}, ``Dog-owner attachment is associated with oxytocin receptor gene polymorphisms in both parties,'' \emph{Frontiers in Psychology}, vol. 9, art. 435, 2018, doi: 10.3389/fpsyg.2018.00435.

\bibitem{rault2020}
J.-L. Rault, S. Waiblinger, X. Boivin, and P. H. Hemsworth, ``The power of a positive human-animal relationship for animal welfare,'' \emph{Frontiers in Veterinary Science}, vol. 7, art. 590867, 2020, doi: 10.3389/fvets.2020.590867.

\bibitem{nagasawa2009}
M. Nagasawa, K. Mogi, and T. Kikusui, ``Attachment between humans and dogs,'' \emph{Japanese Psychological Research}, vol. 51, no. 3, pp. 209--221, 2009, doi: 10.1111/j.1468-5884.2009.00402.x.

\bibitem{topal2005}
J. Topal, M. Gacsi, A. Miklosi, V. Viranyi, Z. Kubinyi, and A. Csanyi, ``Attachment to humans: a comparative study on hand-reared wolves and differently socialized dog puppies,'' \emph{Animal Behaviour}, vol. 70, no. 6, pp. 1367--1375, 2005.

\bibitem{horn2013}
L. Horn, I. Range, and F. Huber, ``The importance of the secure base effect for domestic dogs: evidence from a manipulative problem-solving task,'' \emph{PLOS ONE}, vol. 8, no. 5, e65254, 2013, doi: 10.1371/journal.pone.0065254.

\bibitem{hare2005}
B. Hare and M. Tomasello, ``Human-like social skills in dogs?'' \emph{Trends in Cognitive Sciences}, vol. 9, no. 9, pp. 439--444, 2005.

\bibitem{miklosi2013}
A. Miklosi and J. Topal, ``What does it take for two species to become best friends? Evolutionary changes in canine social competence,'' \emph{Trends in Cognitive Sciences}, vol. 17, no. 6, pp. 287--294, 2013.

\bibitem{udell2016}
M. A. R. Udell and L. Brubaker, ``Are dogs social generalists? Canine social cognition, attachment, and the dog-human bond,'' \emph{Current Directions in Psychological Science}, vol. 25, no. 5, pp. 327--333, 2016.

\bibitem{handlin2015}
L. Handlin \emph{et al.}, ``Short-term interaction between dogs and their owners: effects on oxytocin, cortisol, insulin and heart rate,'' \emph{Physiology \& Behavior}, vol. 145, pp. 103--110, 2015, doi: 10.1016/j.physbeh.2015.03.022.

\bibitem{axelsson2013}
E. Axelsson \emph{et al.}, ``The genomic signature of dog domestication reveals adaptation to a starch-rich diet,'' \emph{Nature}, vol. 495, pp. 360--364, 2013, doi: 10.1038/nature11837.

\bibitem{driscoll2007}
C. A. Driscoll \emph{et al.}, ``The Near Eastern origin of cat domestication,'' \emph{Science}, vol. 317, no. 5837, pp. 519--523, 2007, doi: 10.1126/science.1139518.

\bibitem{ottoni2017}
C. Ottoni \emph{et al.}, ``The palaeogenetics of cat dispersal in the ancient world,'' \emph{Nature Ecology \& Evolution}, vol. 1, art. 0139, 2017, doi: 10.1038/s41559-017-0139.

\bibitem{berns2012}
G. S. Berns, D. M. Brooks, and M. Spivak, ``Functional MRI in awake unrestrained dogs,'' \emph{PLOS ONE}, vol. 7, no. 5, e38027, 2012, doi: 10.1371/journal.pone.0038027.

\bibitem{dietz2018}
L. Dietz, A.-M. K. Arnold, V. C. Goerlich, and C. M. Vinke, ``The importance of early life experiences for the development of behavioural disorders in domestic dogs,'' \emph{Behaviour}, vol. 155, no. 2--3, pp. 83--114, 2018, doi: 10.1163/1568539X-00003486.

\end{thebibliography}

\end{document}
